\documentclass[11pt]{article}
\usepackage[margin=1.1in]{geometry}
\usepackage{amsmath,amssymb,amsthm,mathtools,booktabs,array,tabularx,longtable}
\usepackage[dvipsnames]{xcolor}
\usepackage[colorlinks=true,linkcolor=MidnightBlue,citecolor=MidnightBlue,urlcolor=MidnightBlue]{hyperref}
\usepackage{fancyhdr}
\pagestyle{fancy}\fancyhf{}\renewcommand{\headrulewidth}{0pt}
\fancyfoot[R]{\thepage}
\fancyfoot[L]{\scriptsize\color{gray}Generated by the Transport Lab from approved parts of its store; not a deposited paper. Papers of record: \href{https://zenodo.org/communities/transport-papers}{Zenodo, Transport Papers}}
\fancypagestyle{plain}{\fancyhf{}\fancyfoot[R]{\thepage}\fancyfoot[L]{\scriptsize\color{gray}Generated by the Transport Lab from approved parts of its store; not a deposited paper. Papers of record: \href{https://zenodo.org/communities/transport-papers}{Zenodo, Transport Papers}}\renewcommand{\headrulewidth}{0pt}}
\newtheorem{theorem}{Theorem}[section]
\newtheorem{proposition}[theorem]{Proposition}
\newtheorem{lemma}[theorem]{Lemma}
\newtheorem{corollary}[theorem]{Corollary}
\newtheorem{conjecture}[theorem]{Conjecture}
\newtheorem{refutation}[theorem]{Refutation}
\newtheorem{prediction}[theorem]{Prediction}
\newtheorem{principle}[theorem]{Principle}
\newtheorem{observation}[theorem]{Observation}
\newtheorem{axiom}[theorem]{Axiom}
\theoremstyle{definition}
\newtheorem{definition}[theorem]{Definition}
\newtheorem{fact}[theorem]{Fact}
\newtheorem{assembled}[theorem]{Assembly}
\newtheorem{postulate}[theorem]{Postulate}
\newtheorem{premise}[theorem]{Premise}
\newtheorem{hypothesis}[theorem]{Hypothesis}
\newtheorem{remark}[theorem]{Remark}
\newtheorem{example}[theorem]{Example}
\theoremstyle{plain}
\setlength{\marginparwidth}{0.8in}
\newcommand{\changedmark}[2]{\marginpar{\raggedright\scriptsize\color{#1}#2}}
\providecommand{\Aa}{}\renewcommand{\Aa}{\mathcal{A}}
\providecommand{\Ad}{}\renewcommand{\Ad}{\widetilde{\mathrm{Ad}}}
\providecommand{\Aev}{}\renewcommand{\Aev}{\mathcal{A}^{+}}
\providecommand{\Aod}{}\renewcommand{\Aod}{\mathcal{A}^{-}}
\providecommand{\B}{}\renewcommand{\B}{\mathrm{B}}
\providecommand{\C}{}\renewcommand{\C}{\mathbb{C}}
\providecommand{\CAR}{}\renewcommand{\CAR}{\mathrm{CAR}}
\providecommand{\CS}{}\renewcommand{\CS}{{\mathbb{S}^{+}}}
\providecommand{\CSm}{}\renewcommand{\CSm}{{\mathbb{S}^{-}}}
\providecommand{\Cl}{}\renewcommand{\Cl}{\mathrm{Cl}}
\providecommand{\Con}{}\renewcommand{\Con}{\mathcal{C}}
\providecommand{\Dl}{}\renewcommand{\Dl}{\Delta}
\providecommand{\EN}{}\renewcommand{\EN}{E_{N}}
\providecommand{\End}{}\renewcommand{\End}{\mathrm{End}}
\providecommand{\F}{}\renewcommand{\F}{\mathbb{F}}
\providecommand{\G}{}\renewcommand{\G}{\Gamma}
\providecommand{\HH}{}\renewcommand{\HH}{\mathbb{H}}
\providecommand{\IS}{}\renewcommand{\IS}{I\sigma_3}
\providecommand{\J}{}\renewcommand{\J}{\mathbf{J}}
\providecommand{\LH}{}\renewcommand{\LH}{L_{H}}
\providecommand{\Mat}{}\renewcommand{\Mat}{\mathrm{M}}
\providecommand{\Mo}{}\renewcommand{\Mo}{M^{\circ}}
\providecommand{\Mp}{}\renewcommand{\Mp}{M^{+}}
\providecommand{\Msun}{}\renewcommand{\Msun}{M_{\odot}}
\providecommand{\Mt}{}\renewcommand{\Mt}{\widetilde{M}}
\providecommand{\Ncone}{}\renewcommand{\Ncone}{\mathcal{N}}
\providecommand{\Neff}{}\renewcommand{\Neff}{N_{\mathrm{eff}}}
\providecommand{\PCS}{}\renewcommand{\PCS}{\mathbb{P}{\mathbb{S}^{+}}}
\providecommand{\PT}{}\renewcommand{\PT}{\mathbb{PT}}
\providecommand{\Pin}{}\renewcommand{\Pin}{\mathrm{Pin}}
\providecommand{\Q}{}\renewcommand{\Q}{\mathbb{Q}}
\providecommand{\R}{}\renewcommand{\R}{\mathbb{R}}
\providecommand{\Ric}{}\renewcommand{\Ric}{\mathrm{Ric}}
\providecommand{\SBH}{}\renewcommand{\SBH}{S_{\mathrm{BH}}}
\providecommand{\Sc}{}\renewcommand{\Sc}{\mathbb{S}}
\providecommand{\Sch}{}\renewcommand{\Sch}{\mathrm{P}}
\providecommand{\Sig}{}\renewcommand{\Sig}{\Sigma}
\providecommand{\Som}{}\renewcommand{\Som}{S_\omega}
\providecommand{\Sp}{}\renewcommand{\Sp}{\mathbb{S}}
\providecommand{\Sph}{}\renewcommand{\Sph}{\mathrm{S}}
\providecommand{\Spi}{}\renewcommand{\Spi}{S_\pi}
\providecommand{\Spin}{}\renewcommand{\Spin}{\mathrm{Spin}}
\providecommand{\Sth}{}\renewcommand{\Sth}{S_{\mathrm{th}}}
\providecommand{\T}{}\renewcommand{\T}{\mathbb{T}}
\providecommand{\TH}{}\renewcommand{\TH}{T_{\mathrm{H}}}
\providecommand{\Tr}{}\renewcommand{\Tr}{\operatorname{Tr}}
\providecommand{\UHF}{}\renewcommand{\UHF}{\mathrm{UHF}}
\providecommand{\Z}{}\renewcommand{\Z}{\mathbb{Z}}
\providecommand{\Zc}{}\renewcommand{\Zc}{\mathrm{Z}}
\providecommand{\cmg}{}\renewcommand{\cmg}{\mathrm{cm^{2}\,g^{-1}}}
\providecommand{\dS}{}\renewcommand{\dS}{\mathrm{dS}}
\providecommand{\dsc}{}\renewcommand{\dsc}{\mathrm{dsc}}
\providecommand{\dtPG}{}\renewcommand{\dtPG}{dt}
\providecommand{\dtS}{}\renewcommand{\dtS}{dt_{\mathrm{S}}}
\providecommand{\ellv}{}\renewcommand{\ellv}{\ell}
\providecommand{\er}{}\renewcommand{\er}{e_r}
\providecommand{\et}{}\renewcommand{\et}{e_t}
\providecommand{\ev}{}\renewcommand{\ev}{\mathrm{ev}}
\providecommand{\gl}{}\renewcommand{\gl}{\mathfrak{gl}}
\providecommand{\hatV}{}\renewcommand{\hatV}{\widehat{V}}
\providecommand{\hb}{}\renewcommand{\hb}{\bar{h}}
\providecommand{\hf}{}\renewcommand{\hf}{\hat{f}}
\providecommand{\id}{}\renewcommand{\id}{\mathrm{id}}
\providecommand{\kpc}{}\renewcommand{\kpc}{\mathrm{kpc}}
\providecommand{\lP}{}\renewcommand{\lP}{\ell_{P}}
\providecommand{\mD}{}\renewcommand{\mD}{\mathbf{m}_{D}}
\providecommand{\mb}{}\renewcommand{\mb}{\bar m}
\providecommand{\mnu}{}\renewcommand{\mnu}{\Sigma m_{\nu}}
\providecommand{\nInf}{}\renewcommand{\nInf}{n_\infty}
\providecommand{\nO}{}\renewcommand{\nO}{n_0}
\providecommand{\ninf}{}\renewcommand{\ninf}{n_{\infty}}
\providecommand{\nz}{}\renewcommand{\nz}{n_{0}}
\providecommand{\proj}{}\renewcommand{\proj}{\operatorname{proj}}
\providecommand{\pv}{}\renewcommand{\pv}{\pi}
\providecommand{\rad}{}\renewcommand{\rad}{\mathrm{rad}}
\providecommand{\sgn}{}\renewcommand{\sgn}{\operatorname{sgn}}
\providecommand{\shA}{}\renewcommand{\shA}{\mathcal{A}}
\providecommand{\shC}{}\renewcommand{\shC}{\mathcal{C}}
\providecommand{\shE}{}\renewcommand{\shE}{\mathcal{E}}
\providecommand{\shJ}{}\renewcommand{\shJ}{\mathcal{J}}
\providecommand{\shK}{}\renewcommand{\shK}{\mathcal{K}}
\providecommand{\shL}{}\renewcommand{\shL}{L}
\providecommand{\shW}{}\renewcommand{\shW}{\mathcal{W}}
\providecommand{\sll}{}\renewcommand{\sll}{\mathfrak{sl}}
\providecommand{\so}{}\renewcommand{\so}{\mathfrak{so}}
\providecommand{\su}{}\renewcommand{\su}{\mathfrak{su}}
\providecommand{\supp}{}\renewcommand{\supp}{\operatorname{supp}}
\providecommand{\tr}{}\renewcommand{\tr}{\operatorname{tr}}
\providecommand{\uu}{}\renewcommand{\uu}{\mathfrak{u}}
\providecommand{\br}{}\renewcommand{\br}[1]{\langle #1\rangle}
\providecommand{\bra}{}\renewcommand{\bra}[1]{\langle #1|}
\providecommand{\file}{}\renewcommand{\file}[1]{\texttt{#1}}
\providecommand{\hp}{}\renewcommand{\hp}[2]{( #1 \,|\, #2 )}
\providecommand{\ip}{}\renewcommand{\ip}[2]{\langle #1, #2\rangle}
\providecommand{\ket}{}\renewcommand{\ket}[1]{|#1\rangle}
\providecommand{\rel}{}\renewcommand{\rel}[2]{D\bigl(#1 \,\|\, #2\bigr)}
\providecommand{\rev}{}\renewcommand{\rev}[1]{\widetilde{#1}}
\providecommand{\slashed}{}\renewcommand{\slashed}[1]{{#1}\mkern-10mu/}
\title{\bfseries Neutrino masses}
\author{Leon Butler\\ \normalsize\href{mailto:lbutler.math@gmail.com}{\texttt{lbutler.math@gmail.com}}\\ \normalsize ORCID: \href{https://orcid.org/0009-0006-4449-5617}{0009-0006-4449-5617}}
\date{Generated by the Transport Lab, 2 November 2026}
\begin{document}
\maketitle
\begin{abstract}
\normalsize
Neutrinos are the lightest known matter particles, far lighter than any of their relatives. This trail follows how the framework accounts for them, one step at a time.

It finds a right-handed neutrino, which the Standard Model treats as optional, with its zero hypercharge fixed by the framework's own construction. Assuming a large Majorana mass (Premise (M)), the seesaw makes the light masses small, and because the family couplings start out evenly weighted, only one light neutrino gets a mass at first approximation.

If the second stage of symmetry breaking is hierarchical (Premise (H)), the three masses add up to about 0.059 eV, the smallest total oscillation experiments allow. The walk ends with the matter-antimatter imbalance and with predictions that running experiments can test.
\end{abstract}
\newpage
\noindent \textbf{Authorship and disclosure. }The founding idea of the series is the author's: that division by zero in a Clifford algebra should be read not as a prohibition but as a transport into a higher-dimensional algebra, and that the transported division need not be accompanied by an inverse operation. The direction of investigation and the critical review of the results are likewise the author's. The mathematical development and the drafting of the text were carried out in extended collaboration with Claude Fable 5 (Anthropic), an AI system.

\noindent\fbox{\parbox{\dimexpr\linewidth-2\fboxsep-2\fboxrule}{\textbf{This is not a deposited paper.} It was generated by the Transport Lab on 2 November 2026 and assembled from approved parts of its store. It follows results across several papers of the series; each is stated with its proof and credited to the paper that first proved it. The official versions of those papers are on Zenodo, in the Transport Papers community, \url{https://zenodo.org/communities/transport-papers}.}}
\newpage

\section*{Introduction}
The walk starts with the particle that makes the seesaw possible, sets aside one idea that failed, and then builds the seesaw piece by piece.

It closes with the predictions the seesaw makes and what they say about the early universe.

\section{What we assume}

\begin{postulate}[Hypothesis (U): equal rank-one eigenvalues at the ceiling]\label{c:dark-matter.hyp-u-equal-rank-one-eigenvalues}
This is an assumption, not a result: at the highest energy the framework describes, the basic mass of each kind of matter particle is taken to be one common number. The top quark keeps that value, while smaller later effects pull the others apart. Because the assumption can be read in two ways, the quantities that rest on it come as ranges rather than single values.
\end{postulate}

\begin{postulate}[H: the second stage is hierarchical]\label{c:neutrino.h-second-stage-hierarchical}
This is an assumption, not a result: the lightest neutrino gets almost no mass, less than a hundredth of the next one. Other particle families show steep mass hierarchies of this kind, which makes the assumption plausible without proving it. With it, the framework's total neutrino mass becomes a sharp number rather than a lower limit.
\end{postulate}

\begin{postulate}[M: the framework realises the Majorana term, at a scale $M \gg v$]\label{c:neutrino.m-framework-realises-majorana-term-scale}
Neutrinos could carry a special kind of mass that only electrically neutral particles can have. Nothing forbids it, but nothing in the framework yet shows that it forms or how big it is. The paper assumes that it does, at a very high scale, and leaves the explanation for later work.
\end{postulate}

\begin{postulate}[R: the descent reheats below the lower heavy scale]\label{c:neutrino.r-descent-reheats-below-lower-heavy}
This is an assumption about the early universe: after the event the framework calls the descent, it never became hot enough to produce the heavy neutrinos. If it holds, those neutrinos could neither create nor erase the matter-antimatter imbalance. The paper notes that the assumption is not derived and asks more than its simplest justification would give.
\end{postulate}

\begin{postulate}[The real-breaking assumption]\label{c:neutrino.real-breaking-assumption}
This is an assumption, not a result: the small distortion that gives the lighter neutrinos their masses is taken to involve no complex phase. If it holds, the framework predicts no matter-antimatter difference in how neutrinos mix and two sharp values for a rare nuclear decay rate; if it fails, those sharper predictions are lost.
\end{postulate}

\section{Neutrino masses}

Each step below adds one piece to the account of why neutrinos are so light.

The story starts with a right-handed neutrino, a particle the Standard Model treats as optional. In this framework it is not optional: it is the empty conformal spinor, the state of one conformal spinor with no pair mode occupied.

Its hypercharge is zero, and that value comes from the framework's own construction, since anomaly cancellation alone would have left it free.

\begin{proposition}[The empty conformal spinor (the empty twistor)]\label{c:neutrino.empty-twistor}
The framework's list of particles contains a state that behaves exactly like a right-handed neutrino, a particle the Standard Model treats as optional. The paper also shows that the usual consistency checks on particle charges could not have forced this particle's charge to be zero, since they allow a whole family of values. The zero comes from the framework's own construction, and at that value the particle is invisible to all the consistency checks.
\emph{First proved in} T15, Prop 3.1 (\href{https://doi.org/10.5281/zenodo.22723070}{doi:10.5281/zenodo.22723070}).
\end{proposition}
\noindent\emph{Rests on:} \ref{c:neutrino.imported-ceiling-under-hypothesis-u-arrow}, \ref{c:neutrino.imported-derived-content-singlets}.

With a partner for the neutrino in hand, the next question is what sets the neutrino's mass. One early idea was that the cosmological term of the deformed bracket could do it.

That term lies below the flat term by the factor $\Lambda l^{2}/6$ for a composite of extent $l$, so it never dominates for anything smaller than the horizon. The idea is refuted and kept on the record.

\begin{refutation}[The neutrino mass is set by the cosmological term of the deformed bracket]\label{c:refuted.the-neutrino-mass-is-set-by-the-cosmological-ter}
One idea was that the tiny cosmological constant, entering the framework's mass formula, could set the neutrino mass. The paper shows this effect is far too weak for anything smaller than the observable universe, so the idea is refuted. Neutrino masses come instead from the seesaw, with its heavy scale taken as a premise.
\end{refutation}

What remains is the familiar seesaw, and this step separates what in it is proved from what is assumed. Gauge invariance allows exactly one fermion mass that needs no condensate: a Majorana mass for the right-handed neutrino.

Assuming, as Premise (M), that the framework realises that mass at a large scale $M_{R}$, the leading-order seesaw for one generation gives \[ m\_{\nu} = \frac{m\_{D}\^{}{2}}{M\_{R}}, \] with $m_{D}$ the ordinary Dirac mass, which is why the light mass comes out so small.

\begin{proposition}[The seesaw: what gauge invariance proves, and what is assumed]\label{c:neutrino.seesaw-what-gauge-invariance-proves-what}
The paper first rules out one idea: the cosmological constant's tiny effect on the framework's mass formula is far too weak to give neutrinos their mass. It then recalls that gauge symmetry allows exactly one fermion mass that needs no Higgs field, a Majorana mass for the right-handed neutrino. Assuming, as Premise (M), that the framework realises that mass and makes it large, the familiar seesaw follows at leading order: the light neutrino mass is the square of an ordinary mass divided by the huge one.
\emph{First proved in} T15, Prop 4.1 (\href{https://doi.org/10.5281/zenodo.22723070}{doi:10.5281/zenodo.22723070}).
\end{proposition}
\noindent\emph{Rests on:} \ref{c:mass.exact-join}, \ref{c:neutrino.imported-derived-content-singlets}, \ref{c:neutrino.m-framework-realises-majorana-term-scale}.

The seesaw needs the pattern of neutrino couplings across the three families, so the trail pauses for a piece of pure algebra.

A real symmetric $3 \times 3$ matrix that is unchanged when the families are permuted has only two free numbers: \[ S = \alpha\mathbf{1} + \beta\J, \] with $\J$ the all-ones matrix. It has eigenvalue $\alpha + 3\beta$ on the democratic direction $u' = (1,1,1)$ and $\alpha$ on the plane orthogonal to it, which makes it easy to invert.

\begin{lemma}[$S_{3}$-invariant symmetric matrices]\label{c:neutrino.invariant-symmetric-matrices}
If a three-by-three table of numbers must look the same however the three families are relabelled, it can contain only two independent numbers. That fixes its behaviour completely: one special direction in which all three families are equal, and a plane of directions all treated alike. Inverting such a table is then simple.
\emph{First proved in} T15, Lem 4.3 (\href{https://doi.org/10.5281/zenodo.22723070}{doi:10.5281/zenodo.22723070}).
\end{lemma}

Now the pieces combine. The Dirac couplings treat the three families identically at leading order, so the Dirac mass matrix is \[ \mD = d\,\J, \] with the single Dirac eigenvalue $m_{D} = 3d$ along $u'$.

Under Premise (M) only one light neutrino then gets a mass through the seesaw, and the ordering is normal. If the second stage of the breaking is hierarchical (Premise (H)), the sum sits at the oscillation floor.

\begin{proposition}[The rank-one seesaw: ordering, sum, and the double-beta values]\label{c:neutrino.rank-one-seesaw-ordering-sum-double}
Assuming the framework supplies the heavy neutrino's Majorana mass at a scale far above the electroweak one (Premise (M)), and because the framework's neutrino couplings treat the three families identically at first approximation (a structural-level reading), only one light neutrino gets a mass, through the seesaw. Provided the corrections that break that symmetry are small, that fixes the mass ordering as normal; if the second stage of those corrections is also hierarchical (Premise (H)), the total mass sits at the smallest value oscillation experiments allow, about 0.059 eV, and otherwise that value is only a lower bound. Under the same premises neutrinoless double-beta decay occurs, at rates too small for the next generation of experiments to see: a conditional absence.
\emph{First proved in} T15, Prop 4.4 (\href{https://doi.org/10.5281/zenodo.22723070}{doi:10.5281/zenodo.22723070}).
\end{proposition}
\noindent\emph{Rests on:} \ref{c:neutrino.imported-indifference-faces-democratic-alignment-rank}, \ref{c:neutrino.invariant-symmetric-matrices}, \ref{c:neutrino.seesaw-what-gauge-invariance-proves-what}, \ref{c:neutrino.h-second-stage-hierarchical}, \ref{c:neutrino.m-framework-realises-majorana-term-scale}.

At first approximation only one light neutrino has a mass, so the next step asks how a second one gets its mass.

A Dirac breaking of size $\delta$, written $\mD = d(\J + \delta A)$, lifts the second light mass at order $\delta^{2}M_{3}/\alpha$, with $\alpha$ and $M_{3}$ the two heavy scales. A breaking of the Majorana sector alone lifts nothing.

\begin{lemma}[The breaking relation]\label{c:neutrino.breaking-relation}
At first approximation the framework gives only one light neutrino a mass. This result shows how a small distortion of the neutrino's coupling pattern gives mass to a second one: to leading order the effect grows as the square of the distortion and is amplified by the ratio of two heavy mass scales. Distorting only the heavy neutrinos' masses would not do it. At realistic sizes the approximate formula is only a guide, so the paper computes the masses exactly.
\emph{First proved in} T15, Lem 4.6 (\href{https://doi.org/10.5281/zenodo.22723070}{doi:10.5281/zenodo.22723070}).
\end{lemma}
\noindent\emph{Rests on:} \ref{c:neutrino.invariant-symmetric-matrices}.

This sets up the first prediction that a measurement can check.

Under Premise (H) the neutrino mass sum lies at the normal-ordering floor, \[ \mnu \approx 0.059 \text{ eV} \] Cosmological surveys are now close to this value, so it is a near-term test.

\begin{prediction}[Neutrino mass sum at the oscillation floor]\label{c:pred.neutrino-mass-sum-floor}
If the second stage of the framework's family symmetry breaking is hierarchical (Premise (H)), the framework predicts that the three neutrino masses add up to the smallest total that oscillation experiments allow, about 0.059 eV, with the normal ordering. Cosmological surveys are now close to this value, so it is a near-term test.
\emph{Kill condition:} Below: a physical-prior cosmological bound confirmed below the floor, or confirmed inverted ordering. Above: a sum determined meaningfully above the floor (near 0.075 eV), or m1 measured above a few meV.
\end{prediction}

The same premise bears on neutrinoless double-beta decay, a rare nuclear process that would show neutrinos are their own antiparticles.

Under Premise (H) the effective mass $m_{\beta\beta}$ that sets its rate lies in the hull from $1.3$ to $3.9$ meV, below next-generation sensitivity. A clear signal there would count against the framework, so the predicted absence is itself a test.

\begin{prediction}[No neutrinoless double-beta signal at next-generation sensitivity]\label{c:pred.no-neutrinoless-double-beta}
Neutrinoless double-beta decay is a rare nuclear process that would show neutrinos are their own antiparticles. Assuming the framework's premise (H), that the lightest neutrino is nearly massless, the effective mass that sets the rate comes out at a few thousandths of an electronvolt, between about $1.3$ and $3.9$ meV, with the sharper windows $1.3$ to $1.7$ or $3.5$ to $3.9$ meV if the breaking is also assumed real. That is too small for the next generation of experiments to detect, so a clear signal there would count against the framework. The absence is conditional: it rests on premise (H).
\emph{Kill condition:} A confirmed signal above the predicted window.
\end{prediction}

Heavy neutrinos are also the textbook mechanism of leptogenesis, the standard account of why the universe holds more matter than antimatter.

Here that mechanism must not operate, and it is suppressed in three steps: by the democratic alignment at leading order, by Premise (R) against washout and the resonant pair, and by real breaking for the phases. Premise (R) is an assumption: that the descent reheats below the lower heavy scale, so no heavy neutrino ever thermalised.

\begin{proposition}[Leptogenesis suppressed: the alignment, the reheating premise, and a new stance]\label{c:neutrino.leptogenesis-suppressed-alignment-reheating-premise-new}
Heavy neutrinos are the standard explanation of why the universe has more matter than antimatter, but in this framework that imbalance must be inherited, not made here. The paper shows that at first approximation these heavy neutrinos cannot create an imbalance, that they would neither have erased the inherited one nor made a new one provided the early universe never got hot enough to make them (Premise (R), an assumption), and that if the small symmetry breaking is real no CP-violating phase arises. With the heavy scale set under Hypothesis (U), the same premise also means the universe never got back to the temperature at which the forces unify. That gives a testable stance: the leptonic CP phase should be trivial.
\emph{First proved in} T15, Prop 5.3 (\href{https://doi.org/10.5281/zenodo.22723070}{doi:10.5281/zenodo.22723070}).
\end{proposition}
\noindent\emph{Rests on:} \ref{c:neutrino.imported-ceiling-under-hypothesis-u-arrow}, \ref{c:neutrino.imported-indifference-faces-democratic-alignment-rank}, \ref{c:neutrino.invariant-symmetric-matrices}, \ref{c:neutrino.rank-one-seesaw-ordering-sum-double}, \ref{c:neutrino.r-descent-reheats-below-lower-heavy}, \ref{c:dark-matter.hyp-u-equal-rank-one-eigenvalues}, \ref{c:neutrino.real-breaking-assumption}.

The leptogenesis result ends in a stance that experiment can test.

Under the real-breaking assumption the Dirac CP phase of lepton mixing is trivial, \[ \delta\_{CP} \in \{0, \pi\} \] A measured CP phase away from these two values would refute it.

\begin{prediction}[The leptonic CP phase is trivial]\label{c:pred.leptonic-cp-trivial}
If the small symmetry breaking in the lepton sector is real (the paper's real-breaking assumption), the framework predicts that neutrinos and antineutrinos oscillate in the same way, with no CP violation in lepton mixing. Current long-baseline experiments still allow this if the masses have the normal ordering, which the framework also predicts, but exclude it at three sigma for the inverted ordering, so establishing the inverted ordering would refute this prediction along with that one.
\emph{Kill condition:} Confirmed CP violation in the leptonic sector under normal ordering.
\end{prediction}

\section{Results used from elsewhere}
The proofs above use the following results of the series, stated briefly; each is proved in the paper named.

\begin{fact}[The exact join]\label{c:mass.exact-join}
For two conformal spinors the grip on the frozen point has $m^{2} = 2|B_{\Lambda}|^{2} + \tfrac{\Lambda}{3}C$, with $B_{\Lambda} = \br{\pi^{1}\pi^{2}} + \tfrac{\Lambda}{6}[\omega^{1}\omega^{2}]$.
\emph{Source:} T11, Thm 5.6 (\href{https://doi.org/10.5281/zenodo.22722565}{doi:10.5281/zenodo.22722565}).
\end{fact}

\begin{fact}[Imported: the ceiling under Hypothesis (U), the arrow, and the flag covariance]\label{c:neutrino.imported-ceiling-under-hypothesis-u-arrow}
Imported: the rung scale under Hypothesis (U) and Premise (M), the inherited asymmetry (a declared stance) with the orientation as the framework's own CP source, and the flag covariance of the $\nu_{R}$ identification under Hypothesis (S).
\emph{Source:} T15, Fact 2.5 (\href{https://doi.org/10.5281/zenodo.22723070}{doi:10.5281/zenodo.22723070}).
\end{fact}

\begin{fact}[Imported: the derived content and its singlets]\label{c:neutrino.imported-derived-content-singlets}
The derived generation's singlets are $\nu^{c}, u^{c}, d^{c}, e^{c}$ with $Y = 0, -\tfrac23, \tfrac13, 1$, the hypercharge-zero $\nu^{c}$ at grade zero, while the lepton doublet rests on Hypothesis (S).
\emph{Source:} T15, Fact 2.1 (\href{https://doi.org/10.5281/zenodo.22723070}{doi:10.5281/zenodo.22723070}).
\end{fact}

\begin{fact}[Imported: indifference, the faces, the democratic alignment, and the rank-one texture]\label{c:neutrino.imported-indifference-faces-democratic-alignment-rank}
Under the indifference premise, the conjectural face picture of the families and the Gram-square reading, the leading Yukawa form is democratic, proportional to the all-ones matrix, giving one heavy and two light families at leading order.
\emph{Source:} T15, Fact 2.3 (\href{https://doi.org/10.5281/zenodo.22723070}{doi:10.5281/zenodo.22723070}).
\end{fact}

\section*{Conclusion}
Taken together, the seesaw with democratic Dirac couplings gives one light neutrino its mass first, fixes a normal ordering, and puts the sum at the oscillation floor under Premise (H).

The trail rests on stated assumptions: that the framework realises a large Majorana mass (Premise (M)), that the second stage of breaking is hierarchical (Premise (H)), that the early universe never reheated enough to make heavy neutrinos (Premise (R)), and that the rank-one eigenvalues coincide at the ceiling (Hypothesis (U)). A mass sum away from the floor, a double-beta signal, or a nontrivial leptonic CP phase would each count against the picture.

\end{document}